\section{The full record and its coupled collision and return bridges}
\label{sec:v57-coupled-bridges}
The collision and actual-return bridges are obtained from the same
trajectory. Their joint limit is consequently a graph coupling, not
two independent Gaussian bridges. We prove this joint assertion in a
norm which is total variation in the entire observed record and
bounded-Lipschitz dual in the two paths. This also specifies exactly
which adaptive observations and subsequent selections preserve the
conclusion.

Let $\mathcal C=C([0,1],\R^4)$ and give $\mathcal C^2$ the metric
$d_2((x,y),(x',y'))=\max\{\|x-x'\|_\infty,\|y-y'\|_\infty\}$.
The unit test class $\mathcal F_2$ consists of real functions with
supremum norm and Lipschitz constant at most one. Its dual norm is
$\mathrm{BL}_2^*$. Put
\begin{equation}\label{eq:v57-bridge-graph}
 A_R=c^{-1/2}L_R^{-1},\qquad
 \mathsf H_R=\operatorname{Law}(\mathbb B_{\Omega_R},
                                      A_R\mathbb B_{\Omega_R}).
\end{equation}
The second marginal has covariance $D_R$, by
\eqref{eq:v37-covariance-change}; it is not independent of the first.
All these linear maps are uniformly bounded and continuous in $R$.

On the original source $\nu_R^*$, define
\[
 O_{n,R}=(K_{n,R},N_{n,R},T_{n,R}),\qquad
 Z_{n,R}^{\rm path}=(\mathcal B_{N_{n,R},R},\mathcal Y_{n,R}),
\]
using the paths tied to their own terminals in
\eqref{eq:v41-collision-bridge}--\eqref{eq:v41-return-bridge}.
Let $\mathsf T_{n,R}^z$ be the conditional law of this pair at
$O_{n,R}=z$. A single regular disintegration is used on the standard
Borel spaces. The discrete labels are countable. As in
Section~\ref{sec:v52-path-inversion}, a countable norming class for
$\mathcal F_2$ fixes all roof versions off one null set for each set of
labels. Tests measurable in $z$ are interpreted through this same kernel.

\subsection{The joint, rather than marginal, local numerator}
\begin{lemma}[A coupled path law on central roof windows]
\label{lem:v57-joint-local}
For each fixed $M,h<\infty$, $h>0$, uniformly in the original exact
labels with $|Z(t)|\le M$,
\begin{equation}\label{eq:v57-joint-local}
 \frac1h\int_t^{t+h}
 \left\|m^2p_{n,R}(k,m,u)\mathsf T_{n,R}^{k,m,u}
          -G_{m,R}^{n,k}(u)\mathsf H_R\right\|_{\mathrm{BL}_2^*}
                                      du\longrightarrow0
                                      \quad(m\to\infty).
\end{equation}
There is no positive-denominator assumption.
\end{lemma}
\begin{proof}
Under the finite positive source
$\lambda_m=(m^2/h)\one_{\{K_n=k,N_n=m,T_n\in[t,t+h]\}}\nu_R^*$,
Lemma~\ref{lem:v52-return-transfer} proves the actual clock coupling
\[
 \int\min\{2,\|\mathcal Y_{n,R}
       -A_{m,n,R}\mathcal B_{m,R}\|_\infty\}\,d\lambda_m\to0,
 \qquad A_{m,n,R}=\sqrt{m/n}\,L_R^{-1}.
\]
This is \eqref{eq:v52-clock-coupling} on the joint trajectory, not
an inference from convergence of two marginal laws. On the stated
central targets, $n/m=c+O_M(m^{-1/2})$ and hence
$A_{m,n,R}\to A_R$ uniformly. The collision-path measures under
$\lambda_m$ are uniformly tight and have bounded mass, as proved
there. Splitting at $\|\mathcal B_{m,R}\|_\infty\le K$, then taking
$m\to\infty$ followed by $K\to\infty$, permits $A_{m,n,R}$ to be
replaced by $A_R$ in the last display. This argument also covers a
source whose total mass tends to zero.

The graph map $x\mapsto(x,A_Rx)$ is uniformly Lipschitz. Pushing the
collision path numerator of Theorem~\ref{thm:v52-path-remainder}
through this map costs at most $\max\{1,\|A_R\|\}$ in its dual norm.
For every test in $\mathcal F_2$, the difference between the original
pair and this graph is bounded by the truncated clock discrepancy
above. Joint disintegration and the countable norming class allow the
supremum over path tests before roof integration. These two bounds
prove \eqref{eq:v57-joint-local}. The roof interval enlarges the
central target only to $M+h/\sqrt m$, as in the inherited clock lemma.
\end{proof}

\subsection{No label or roof cutoff in the joint limit}
For two finite measures on $\mathcal O_n\times\mathcal C^2$ define
\[
 \|\sigma\|_{\mathrm{obs},\mathrm{BL}}
  =\sup_\phi\left|\int\phi(z,x,y)\,d\sigma\right|,
\]
where $\phi$ is jointly measurable and $\phi(z,\cdot,\cdot)\in\mathcal F_2$
for every $z$. For a disintegrated signed measure this equals the
integral of the fiber dual norms: measurable selection of a test from
the countable norming class gives the lower inequality, and integration
gives the upper one. This norm is variation mass on the observation
marginal, not probability total variation; it is not path-space total
variation.

\begin{theorem}[Untruncated mixed law with coupled bridges]
\label{thm:v57-global-coupled}
Uniformly in $R$ at each fixed actual return index $n$,
\begin{equation}\label{eq:v57-global-pair-error}
 \epsilon_n:=\sup_R\int_{\mathcal O_n}
 \left\|p_{n,R}(z)\mathsf T_{n,R}^{z}
                  -\gamma_{n,R}(z)\mathsf H_R\right\|_{\mathrm{BL}_2^*}
                          d\zeta\longrightarrow0\quad(n\to\infty).
\end{equation}
Consequently, with the normalized positive arithmetic law
$\mathsf A_{n,R}$ of \eqref{eq:v57-full-reference},
\begin{equation}\label{eq:v57-global-joint-law}
 \Delta_n:=\sup_R
 \left\|\operatorname{Law}_{\nu_R^*}(O_{n,R},Z_{n,R}^{\rm path})
                   -\mathsf A_{n,R}\otimes\mathsf H_R
       \right\|_{\mathrm{obs},\mathrm{BL}}
                          \le2\epsilon_n\longrightarrow0.
\end{equation}
In particular the conditional pair bridge at its own complete observed
record satisfies
\begin{equation}\label{eq:v57-global-conditional-mean}
 \sup_R\int d_{\mathrm{BL}_2}(\mathsf T_{n,R}^z,\mathsf H_R)
                                             \,d\mathsf P_{n,R}(z)
                                             \le2\epsilon_n\to0.
\end{equation}
All three statements are on the original unprotected source.
\end{theorem}
\begin{proof}
On $|V_{n,R}|\le M$, the discrete-label and unit-window count is the
same $O_M(n^2)$ count as in \eqref{eq:v57-central-error}. The collision
central coordinates are bounded by a constant depending on $M$,
since $m\asymp n$. Thus Lemma~\ref{lem:v57-joint-local}, including the
fixed enlargement across each unit window, makes the integral of the
joint dual error on this region tend uniformly to zero.

Outside that region its integrand is at most
$p_{n,R}+|\gamma_{n,R}|$: positive finite path measures have dual
norm equal to mass. The true fourth-moment tail and
\eqref{eq:v57-reference-tail} bound its total by
$CM^{-4}+Ce^{-aM^2}+Ce^{-an}$. First take the return limit and then
the central radius to infinity. This proves
\eqref{eq:v57-global-pair-error} on the entire observation space.

For $p\ge0$, a probability $Q$ and a probability $H$, replacing a
negative scalar $g$ by zero decreases
$\|pQ-gH\|_{\mathrm{BL}_2^*}$: the old measure is then positive
with mass $p+|g|$, and the new measure has mass $p$. Thus positive-part
replacement never enlarges the error in
\eqref{eq:v57-global-pair-error}. The constant test one gives
$|Z_{n,R}^{\rm ref}-1|\le\epsilon_n$. Normalization adds exactly
that mass difference in the mixed norm, proving
\eqref{eq:v57-global-joint-law}.

Finally, pointwise in the common disintegration,
\[
 p\|Q-H\|_{\mathrm{BL}_2^*}
 \le\|pQ-gH\|_{\mathrm{BL}_2^*}+|p-g|
 \le2\|pQ-gH\|_{\mathrm{BL}_2^*}.
\]
Integrate with $g=\gamma_{n,R}$ to obtain
\eqref{eq:v57-global-conditional-mean}. Neither a pointwise floor
nor division by a zero fiber mass has been used; an arbitrary version
on $p=0$ has zero weight in this integral.
\end{proof}

\subsection{Observations, randomization and the exact selection event}
A Markov kernel $\mathscr K_{n,R}(z,do)$ from $\mathcal O_n$ to a
standard Borel output space represents any observation of the complete
record with independent randomization. It may, for example, output a
label, a roof interval, an acceptance decision and the random seed.
The output law is formed with this very kernel under both the true
law and the arithmetic reference. A procedure observing the path
itself is not a kernel of the record alone.

\begin{theorem}[Observation and postselection with a mass budget]
\label{thm:v57-selection}
Apply any such kernel to the two probability measures in
\eqref{eq:v57-global-joint-law}, leaving the paths unchanged. Their
output--path mixed norm is at most $\Delta_n$, uniformly over all
these kernels. Let $A$ be a measurable output event and suppose its
reference probability is at least $r_n>\Delta_n$. Then its true
probability is at least $r_n-\Delta_n$, and the conditional laws obey
\begin{align}
 d_{\rm TV}(\text{true output}\mid A,
                       \text{reference output}\mid A)
        &\le\frac{\Delta_n}{r_n-\Delta_n},
                                      \label{eq:v57-selection-TV}\\
 \|\text{true output--path law}\mid A
       -\text{reference output law}\mid A\otimes\mathsf H_R
       \|_{\mathrm{obs},\mathrm{BL}}
        &\le\frac{2\Delta_n}{r_n-\Delta_n}.
                                      \label{eq:v57-selection-joint}
\end{align}
The true conditional mean of the same-output pair-bridge distance
from $\mathsf H_R$ is bounded by the right side of
\eqref{eq:v57-selection-joint}. Hence all these errors vanish for
any selection family with $\Delta_n/r_n\to0$.
\end{theorem}
\begin{proof}
An admissible output--path test pulls back to
$\int\phi(o,x,y)\mathscr K_{n,R}(z,do)$, which is still bounded by
one and one-Lipschitz in the two paths, measurably in $z$. This proves
contraction of the mixed norm, without any regularity of the
observation kernel in the record coordinates. Under the reference,
the graph bridge remains independent of the output because the kernel
does not read the paths.

Write $\alpha,\beta$ for the true and reference probabilities of $A$.
Their difference is at most $\Delta_n$, since the indicator of $A$
is an admissible constant path test. Restrict both joint laws to $A$.
Their mixed-norm error is at most $\Delta_n$. Adding and subtracting
the reference restriction divided by $\alpha$ shows that normalization
costs at most
$(\Delta_n+|\alpha-\beta|)/\alpha\le2\Delta_n/(r_n-\Delta_n)$.
On output marginals this is an $L^1$ variation bound, so division by
two gives \eqref{eq:v57-selection-TV}.

Disintegrate the output--path law. At each output use the last
pointwise inequality in the proof of
Theorem~\ref{thm:v57-global-coupled}, now with the reference output
density as $g$ with respect to a common dominating output measure.
Its integral on $A$ is at most $2\Delta_n$. Division by
$\alpha\ge r_n-\Delta_n$ proves the conditional-mean assertion.
The reference output measure can be discrete, continuous or mixed;
no roof representative is needed in this argument.
\end{proof}

\paragraph{Interpretation of selection.}
The event in Theorem~\ref{thm:v57-selection} is the full output
selection event. For an interval chosen from the observed record,
conditioning on its reported value is not silently replaced by
conditioning only on membership in that interval. Both laws use the
same kernel, including all information carried by selection. Arbitrary
path-dependent acceptance need not preserve a Gaussian bridge
amplitude and is not asserted to do so. Nor does the qualitative
sequence $\Delta_n$ provide a rate for a prescribed microscopic
singleton whose mass is much smaller than $\Delta_n$.

\paragraph{Proof route and the original endpoint.}
The route to \eqref{eq:v57-global-joint-law} consists of four inputs:
centered positive pressure and the physical visible-branch pairing;
the resulting uniform Gaussian envelope for the full transition sum;
the inherited complete-source scalar and path local-variation laws;
and the actual joint clock coupling. The scalar fourth moment pays
the discarded tails in the proof. All cutoffs disappear from the final
probability measures. The two paths remain coupled, with the second
one obtained from the first by $A_R$ in the limit.

This integrated mixed norm does not control a density spike at an
individual roof, even though it controls all bounded observations of
the full record. The incidence and clearance essential-height
conditions and unrestricted pointwise bridges remain the same targets
as in Corollary~\ref{cor:v51-positive-criterion}. The arithmetic
transition kernel is permanent in the uniform theorem; an unmodulated
specialization requires the separate concrete residue criterion.
